\documentclass[../main.tex]{subfiles}

\begin{document}
\section{Results and Discussion}\label{section.4}

\subsection{Hyper-parameter Tuning Results}
\noindent
Tables~\ref{tab.svm-tune}--\ref{tab.cnn-tune} list the top ten
hyper-parameter combinations for each model, ranked by validation accuracy
on the held-out 6{,}400-sample partition (or by mean three-fold
cross-validated accuracy on the 5{,}000-sample subset for the SVM).

\begin{table}[H]
    \centering
    \caption{SVM grid search — top 10 configurations on the 5{,}000-sample
    tuning subset (3-fold CV).}
    \label{tab.svm-tune}
    \begin{tabular}{ccccc}
        \toprule
        \textbf{PCA $d$} & \textbf{$C$} & \textbf{$\gamma$} &
        \textbf{Mean fit time (s)} & \textbf{CV accuracy (\%)} \\
        \midrule
        30 & 5.0 & 0.001 & 0.65 & 63.20 \\
        30 & 5.0 & scale & 0.56 & 63.02 \\
        30 & 1.0 & 0.001 & 0.60 & 62.26 \\
        50 & 5.0 & 0.001 & 0.89 & 62.14 \\
        50 & 5.0 & scale & 0.76 & 62.08 \\
        30 & 1.0 & scale & 0.60 & 61.50 \\
        50 & 1.0 & 0.001 & 0.82 & 61.32 \\
        30 & 0.5 & 0.001 & 0.74 & 61.08 \\
        50 & 1.0 & scale & 0.76 & 61.02 \\
        100 & 1.0 & 0.001 & 1.37 & 60.70 \\
        \bottomrule
    \end{tabular}
\end{table}

\begin{table}[H]
    \centering
    \caption{MLP grid search — top 10 configurations on the validation
    partition.}
    \label{tab.mlp-tune}
    \begin{tabular}{cccccc}
        \toprule
        \textbf{Units} & \textbf{Dropout} & \textbf{LR} &
        \textbf{Epochs} & \textbf{Fit (s)} &
        \textbf{Val acc.\ (\%)} \\
        \midrule
        256 & 0.0 & 1e-03 & 24 & 11.5 & 52.50 \\
        64 & 0.0 & 1e-03 & 19 & 4.6 & 51.59 \\
        256 & 0.0 & 1e-04 & 25 & 11.5 & 51.13 \\
        128 & 0.0 & 1e-04 & 25 & 7.7 & 49.56 \\
        256 & 0.3 & 1e-04 & 25 & 12.4 & 49.02 \\
        64 & 0.0 & 1e-04 & 25 & 5.6 & 48.03 \\
        128 & 0.0 & 1e-03 & 8 & 2.9 & 46.94 \\
        256 & 0.5 & 1e-04 & 25 & 12.2 & 46.89 \\
        64 & 0.3 & 1e-04 & 25 & 5.8 & 46.67 \\
        128 & 0.3 & 1e-04 & 25 & 8.0 & 46.09 \\
        \bottomrule
    \end{tabular}
\end{table}

\begin{table}[H]
    \centering
    \caption{CNN grid search — top 10 configurations on the validation
    partition.}
    \label{tab.cnn-tune}
    \begin{tabular}{ccccccc}
        \toprule
        \textbf{Filters} & \textbf{Kernel} & \textbf{Dropout} &
        \textbf{LR} & \textbf{Epochs} & \textbf{Fit (s)} &
        \textbf{Val acc.\ (\%)} \\
        \midrule
        16 & 5 & 0.3 & 1e-03 & 25 & 84.1 & 75.17 \\
        16 & 3 & 0.3 & 1e-03 & 23 & 65.9 & 74.97 \\
        32 & 3 & 0.3 & 1e-03 & 18 & 61.4 & 73.53 \\
        32 & 5 & 0.3 & 1e-03 & 14 & 69.1 & 72.84 \\
        32 & 3 & 0.0 & 1e-04 & 25 & 92.8 & 71.20 \\
        32 & 3 & 0.3 & 1e-04 & 25 & 91.4 & 70.53 \\
        32 & 3 & 0.0 & 1e-03 & 11 & 38.0 & 70.34 \\
        32 & 5 & 0.3 & 1e-04 & 25 & 113.9 & 70.20 \\
        32 & 5 & 0.0 & 1e-03 & 9 & 39.5 & 67.47 \\
        16 & 3 & 0.0 & 1e-04 & 25 & 65.4 & 67.27 \\
        \bottomrule
    \end{tabular}
\end{table}

\noindent
Per-axis bar charts of mean and best validation accuracy
(Appendix~\ref{appendix.tuning}, Figures~\ref{fig.tune-svm}--\ref{fig.tune-cnn})
make the trends inside each grid easier to read at a glance and are
referenced from the discussion below.

\subsection{Hyper-parameter Tuning Discussion}

\subsubsection{SVM}
\noindent
The best SVM configuration on the tuning subset uses
$d{=}30$ \gls{PCA} components,
$C{=}5.0$ and
$\gamma{=}0.001$, reaching
$63.20\%$ mean cross-validated accuracy.
Two trends are clear from Table~\ref{tab.svm-tune}: the kernel bandwidth
$\gamma$ has by far the largest effect — the worst combinations all use
$\gamma{=}0.01$, which produces an overly localised RBF that
under-fits the 5{,}000-sample subset — and increasing the number of
\gls{PCA} components past $\sim$\,30 does not help on this dataset. The
latter is consistent with a steep PathMNIST principal-component spectrum:
the first few dozen directions already explain most of the
class-relevant variance, while later components mostly model noise. The
regularisation strength $C$ contributes a smaller, mostly monotone
improvement up to $C{=}5$ for the better $\gamma$ values.

\subsubsection{MLP}
\noindent
The MLP grid shows that the learning rate dominates the outcome and that
the dense network still performs substantially worse than the CNN despite
moderate improvements from tuning. The top configuration uses
units${=}256$,
dropout${=}0.0$ and
$\eta{=}10^{-3}$,
reaching approximately
$52.5\%$ on the validation partition. Two trends are visible. First, $\eta{=}10^{-3}$ consistently outperforms $\eta{=}10^{-4}$,
suggesting that the MLP benefits from faster optimisation dynamics on
this relatively small low-resolution dataset.
Second, dropout in this regime is mostly a hindrance — $p{=}0$ wins at
every width because there is barely enough signal in the dense
representation for further regularisation to help. Width has a moderate positive effect overall, with the best-performing
configuration using 256 hidden units. However, even with increased capacity, the MLP remains clearly below the
performance of the convolutional approach.

\subsubsection{CNN}
\noindent
Two effects dominate the CNN grid. First, the larger learning rate
$\eta{=}10^{-3}$ is clearly preferred: every configuration in the
top four uses it, and the deeper feature hierarchy absorbs the faster
optimiser without diverging. Second, moderate dropout ($p{=}0.3$)
helps consistently — every top-four configuration uses it and the next
band of accuracies all uses $p{=}0$, so dropout buys roughly two
percentage points on the validation partition. The choice of filter
count and kernel size matters less than expected at this resolution:
$f{=}16$ with
$k{=}5$ wins narrowly, but
$f{=}32$ and $k{=}3$ are within roughly two percentage points. Smaller filter
counts also train substantially faster, which is one reason we keep
$f{=}16$ for the final model.

\subsection{Final Results}
\noindent
Each model was retrained with its best hyper-parameters on the full
32{,}000-sample training set and evaluated on the 8{,}000-sample
PathMNIST test set. Table~\ref{tab.final} summarises the result.

\begin{table}[H]
    \centering
    \caption{Final-model performance and runtime on the PathMNIST test
    set.}
    \label{tab.final}
    \begin{tabularx}{\linewidth}{lcccc>{\raggedright\arraybackslash}X}
        \toprule
        \textbf{Model} & \textbf{Acc.\ (\%)} &
        \textbf{Macro $F_{1}$ (\%)} & \textbf{Train (s)} &
        \textbf{Predict (s)} & \textbf{Hyper-parameters} \\
        \midrule
        SVM & 69.01 & 68.26 & 10.7 & 7.0 & $d{=}30$, $C{=}5.0$, $\gamma{=}0.001$ \\
        MLP & 56.90 & 53.50 & 15.3 & 0.16 & units${=}256$, dropout${=}0.0$, $\eta{=}10^{-3}$ \\
        CNN & 76.76 & 75.63 & 65.0 & 0.6 & filters${=}16$, $k{=}5$, dropout${=}0.3$, $\eta{=}10^{-3}$ \\
        \bottomrule
    \end{tabularx}
\end{table}

\subsection{Results Discussion and Analysis}
\noindent
On the test set the CNN achieves the highest accuracy
(76.76\%), clearly outperforming both the SVM (69.01\%) and the
MLP (56.90\%). The CNN improves on the SVM by almost
8 percentage points, while the gap between the CNN and MLP
is approximately 20 percentage points.

\noindent
The relative ordering of the three classifiers strongly matches the
theoretical expectations discussed in
Section~\ref{section.3}. The \gls{MLP} performs substantially worse than the CNN it treats
every pixel independently and therefore cannot efficiently exploit the
local spatial structure and texture continuity that dominate
histopathology images. Although the \gls{SVM} benefits from the smooth
non-linear decision boundary of the \gls{RBF} kernel and performs
reasonably well after dimensionality reduction, it still operates on
flattened image vectors and therefore loses important spatial
relationships between neighbouring pixels.

\noindent
The \gls{CNN} clearly achieves the strongest performance because its
convolutional structure directly encodes locality and translation-aware
feature extraction. These inductive biases are particularly important on
PathMNIST, where tissue discrimination depends heavily on subtle local
texture patterns, stain gradients and cellular organisation. Even the
relatively shallow two-block \gls{CNN} used in this assignment is able to
learn substantially more effective representations than either the
kernel-based or fully connected alternatives.

\noindent
The macro-$F_{1}$ values track accuracy closely
(68.26\%, 53.50\% and 75.63\% for the SVM, MLP
and CNN respectively), confirming that the mild class imbalance has not
skewed performance towards the majority class. The row-normalised
confusion matrices in Figure~\ref{fig.cm} show two stable failure modes
shared by all three classifiers: \emph{Cancer-associated Stroma} versus
\emph{Smooth Muscle}, and \emph{Debris} versus \emph{Mucus} — both pairs
share a similar texture and stain profile, exactly as predicted from the
qualitative grid in Figure~\ref{fig.samples}. The MLP additionally shows broader confusion across several visually
similar tissue classes, particularly involving
\emph{Colorectal Adenocarcinoma Epithelium}, which explains the larger
gap between its accuracy and macro-$F_{1}$.

\begin{figure}[!htbp]
    \centering
    \begin{subfigure}[t]{0.32\linewidth}
        \includegraphics[width=\linewidth]{Photos/cm_svm.png}
        \caption{SVM (RBF kernel).}
        \label{fig.cm-svm}
    \end{subfigure}\hfill
    \begin{subfigure}[t]{0.32\linewidth}
        \includegraphics[width=\linewidth]{Photos/cm_mlp.png}
        \caption{MLP.}
        \label{fig.cm-mlp}
    \end{subfigure}\hfill
    \begin{subfigure}[t]{0.32\linewidth}
        \includegraphics[width=\linewidth]{Photos/cm_cnn.png}
        \caption{CNN.}
        \label{fig.cm-cnn}
    \end{subfigure}
    \caption{Row-normalised test-set confusion matrices for the final
    models. The \emph{Debris}/\emph{Mucus} and
    \emph{Cancer-associated Stroma}/\emph{Smooth Muscle} confusions are most
    pronounced in the MLP and SVM models, while the CNN reduces these errors
    substantially through spatial feature extraction. The MLP additionally
    shows broader confusion across several visually similar tissue classes,
    particularly involving \emph{Colorectal Adenocarcinoma Epithelium}.}
    \label{fig.cm}
\end{figure}

\noindent
Runtime highlights a clear trade-off between predictive performance and
computational cost. Training the \gls{SVM} on the full
32{,}000-sample training set took
$10.7\,\text{s}$ on a single CPU thread, while the
\gls{MLP} required
$15.3\,\text{s}$ and the
\gls{CNN}
$65.0\,\text{s}$ on the Apple~M3 \gls{GPU} via
\texttt{tensorflow-metal}~\cite{tfmetal2024}. The substantially longer
training time of the \gls{CNN} reflects the repeated gradient-based
optimisation of multiple convolutional layers and feature maps, whereas
the \gls{SVM} benefits from aggressive dimensionality reduction through
\gls{PCA}.

\noindent
Inference is where the kernel \gls{SVM} pays its theoretical cost:
scoring the 8{,}000 test images takes
$6.95\,\text{s}$ because every test point must be compared against the
learned support vectors, while both neural networks predict in well
under a second ($0.16\,\text{s}$ for the \gls{MLP} and
$0.60\,\text{s}$ for the \gls{CNN}). Although the
\gls{MLP} achieves the fastest inference time overall, its classification
performance remains substantially weaker because it lacks the spatial
inductive biases required for histopathology image analysis.

\noindent
Overall, the \gls{CNN} provides the strongest balance for this dataset.
While its training cost is the highest, it achieves substantially better
accuracy and macro-$F_{1}$ performance than both alternative classifiers
while still maintaining very fast inference. The \gls{SVM} remains an
attractive lightweight baseline because of its relatively low training
cost and competitive accuracy, but its slower inference and inability to
fully exploit spatial structure ultimately limit its scalability for
real-time deployment scenarios.

\end{document}
